% --- Archon physics formalization source begin ---
% archon:physics
% archon:covers PhyXMiniProblems/problem_phyx_mini_0503.lean
% archon:source-report reports/phyx_mini/problem_phyx_mini_0503.source.json

\chapter{Physics problem phyx\_mini\_0503}
\label{ch:PhyXMiniProblems_problem_phyx_mini_0503}

\paragraph{Problem source.}
\#\# Physical scenario

Figure is a graph of \textbackslash{}( |\textbackslash{}psi(x)|\textasciicircum{}2 \textbackslash{}) for a neutron.

\#\# Question

What is the probability that the neutron is located at a posi- tion with \textbackslash{}( |x| \textbackslash{}ge 2\textbackslash{}text\{ fm\} \textbackslash{})?

\#\# Answer choices

- A: A: 0.65

- B: B: 0.55

- C: C: 0.75

- D: D: 0.50

\#\# Figure caption (auxiliary; use the image as primary evidence)

The image is a graph depicting a mathematical function, specifically \textbackslash{}(|\textbackslash{}psi(x)|\textasciicircum{}2\textbackslash{}), plotted against the x-axis labeled "x (fm)," which likely refers to position in femtometers. Key features include:

1. **Axes:**
   - The horizontal axis (x-axis) ranges from -4 to 4 with increments at -4, -2, 0, 2, and 4.
   - The vertical axis represents the function \textbackslash{}(|\textbackslash{}psi(x)|\textasciicircum{}2\textbackslash{}).

2. **Graph:**
   - Two symmetrical linear lines form a "V" shape, intersecting at the origin (0,0).
   - The lines reach the x values of -4 and 4 on the horizontal axis.

3. **Labels:**
   - The vertical axis is labeled with \textbackslash{}(|\textbackslash{}psi(x)|\textasciicircum{}2\textbackslash{}).
   - The letter "a" marks a vertical position above the vertex, on the vertical line.

This likely represents a probability density function in a quantum mechanics context, given the notation \textbackslash{}(|\textbackslash{}psi(x)|\textasciicircum{}2\textbackslash{}).

\#\# Dataset metadata

Quantum Phenomena; Physical Model Grounding Reasoning, Numerical Reasoning

\paragraph{Recorded answer/context.}
C: C: 0.75

\paragraph{Figure/image path.}
/root/proposal\_for\_physic/science-mango/phyx\_mini\_run/phyx\_data/test\_image/503.png

\paragraph{Formalization target.}
create a compiling Lean file with sorry bodies at `PhyXMiniProblems/problem\_phyx\_mini\_0503.lean`. The Lean declarations must preserve the physical quantities, dimensions or dimensional roles, figure labels, governing-law hypotheses, and final relation expressed by this problem.
Use Mathlib/Physlib names found through LeanExplore where available. If a physics API is missing, introduce faithful local abstractions rather than scalar placeholder aliases.

\begin{theorem}[Physics formalization target]
\label{thm:physics:phyx_mini_0503:target}
\lean{PhyXMiniProblems.ProblemPhyXMini0503.problem_phyx_mini_0503}
\uses{def:physics:phyx-mini-0503:phyxminiproblems-problemphyxmini0503-onedimensionalquantumpositionsetup, def:physics:phyx-mini-0503:phyxminiproblems-problemphyxmini0503-requestedpositionprobability, def:physics:phyx-mini-0503:phyxminiproblems-problemphyxmini0503-matchesneutronpositionscenario, def:physics:phyx-mini-0503:phyxminiproblems-problemphyxmini0503-matchesproblemandfigurereadouts, def:physics:phyx-mini-0503:phyxminiproblems-problemphyxmini0503-matchessuppliedprobabilitydensityfigure, def:physics:phyx-mini-0503:phyxminiproblems-problemphyxmini0503-hasphysicalprobabilitydensityparameters, def:physics:phyx-mini-0503:phyxminiproblems-problemphyxmini0503-satisfiesonedimensionalbornrule, def:physics:phyx-mini-0503:phyxminiproblems-problemphyxmini0503-isuniquematchinganswerchoice}
The assigned autoformalize agent should translate this physics problem into Lean declarations in the covered file, with theorem and lemma proof bodies written as `by sorry`.
\end{theorem}
\begin{proof}
This is an autoformalization task, not a proof task. Produce statements that can later be proved by the physics prover without weakening the physical model.
\end{proof}

% --- Archon declaration topology begin ---
\section{Declaration topology}
\begin{definition}[Particle Species]
  \label{def:physics:phyx-mini-0503:phyxminiproblems-problemphyxmini0503-particlespecies}
  \lean{PhyXMiniProblems.ProblemPhyXMini0503.ParticleSpecies}
  Particle species distinguished by the physical scenario
\end{definition}
\begin{proof}
  This is the declared model definition of Particle Species.
\end{proof}
\begin{definition}[Figure Label]
  \label{def:physics:phyx-mini-0503:phyxminiproblems-problemphyxmini0503-figurelabel}
  \lean{PhyXMiniProblems.ProblemPhyXMini0503.FigureLabel}
  Literal labels printed on the supplied probability-density graph
\end{definition}
\begin{proof}
  This is the declared model definition of Figure Label.
\end{proof}
\begin{definition}[Horizontal Tick]
  \label{def:physics:phyx-mini-0503:phyxminiproblems-problemphyxmini0503-horizontaltick}
  \lean{PhyXMiniProblems.ProblemPhyXMini0503.HorizontalTick}
  The five signed position-coordinate ticks visible on the horizontal axis
\end{definition}
\begin{proof}
  This is the declared model definition of Horizontal Tick.
\end{proof}
\begin{definition}[Supplied Probability Density Figure]
  \label{def:physics:phyx-mini-0503:phyxminiproblems-problemphyxmini0503-suppliedprobabilitydensityfigure}
  \lean{PhyXMiniProblems.ProblemPhyXMini0503.SuppliedProbabilityDensityFigure}
  \uses{def:physics:phyx-mini-0503:phyxminiproblems-problemphyxmini0503-figurelabel, def:physics:phyx-mini-0503:phyxminiproblems-problemphyxmini0503-horizontaltick}
  Independent data retained from the primary raster. The plotted values have the dimensional role of inverse-femtometer density readouts. No event probability or answer choice is stored in the figure
\end{definition}
\begin{proof}
  This is the declared model definition of Supplied Probability Density Figure.
\end{proof}
\begin{definition}[One Dimensional Quantum Position Setup]
  \label{def:physics:phyx-mini-0503:phyxminiproblems-problemphyxmini0503-onedimensionalquantumpositionsetup}
  \lean{PhyXMiniProblems.ProblemPhyXMini0503.OneDimensionalQuantumPositionSetup}
  \uses{def:physics:phyx-mini-0503:phyxminiproblems-problemphyxmini0503-particlespecies, def:physics:phyx-mini-0503:phyxminiproblems-problemphyxmini0503-suppliedprobabilitydensityfigure}
  A one-dimensional position model for the neutron. The complex function is the pointwise wave-function amplitude in the selected coordinate unit. Its squared norm is a density readout, while positionProbability is normalized by its ProbabilityMeasure type
\end{definition}
\begin{proof}
  This is the declared model definition of One Dimensional Quantum Position Setup.
\end{proof}
\begin{definition}[position Density Readout]
  \label{def:physics:phyx-mini-0503:phyxminiproblems-problemphyxmini0503-positiondensityreadout}
  \lean{PhyXMiniProblems.ProblemPhyXMini0503.positionDensityReadout}
  \uses{def:physics:phyx-mini-0503:phyxminiproblems-problemphyxmini0503-onedimensionalquantumpositionsetup}
  The probability-density readout |ψ(x)|² at a signed coordinate in the selected length unit
\end{definition}
\begin{proof}
  This is the declared model definition of position Density Readout.
\end{proof}
\begin{definition}[outside Central Region]
  \label{def:physics:phyx-mini-0503:phyxminiproblems-problemphyxmini0503-outsidecentralregion}
  \lean{PhyXMiniProblems.ProblemPhyXMini0503.outsideCentralRegion}
  Positions whose absolute coordinate is at least the supplied cutoff
\end{definition}
\begin{proof}
  This is the declared model definition of outside Central Region.
\end{proof}
\begin{definition}[requested Position Event]
  \label{def:physics:phyx-mini-0503:phyxminiproblems-problemphyxmini0503-requestedpositionevent}
  \lean{PhyXMiniProblems.ProblemPhyXMini0503.requestedPositionEvent}
  \uses{def:physics:phyx-mini-0503:phyxminiproblems-problemphyxmini0503-onedimensionalquantumpositionsetup, def:physics:phyx-mini-0503:phyxminiproblems-problemphyxmini0503-outsidecentralregion}
  The event whose probability is requested in the problem
\end{definition}
\begin{proof}
  This is the declared model definition of requested Position Event.
\end{proof}
\begin{definition}[requested Position Probability]
  \label{def:physics:phyx-mini-0503:phyxminiproblems-problemphyxmini0503-requestedpositionprobability}
  \lean{PhyXMiniProblems.ProblemPhyXMini0503.requestedPositionProbability}
  \uses{def:physics:phyx-mini-0503:phyxminiproblems-problemphyxmini0503-onedimensionalquantumpositionsetup, def:physics:phyx-mini-0503:phyxminiproblems-problemphyxmini0503-requestedpositionevent}
  The requested dimensionless probability, expressed as a real readout
\end{definition}
\begin{proof}
  This is the declared model definition of requested Position Probability.
\end{proof}
\begin{definition}[Matches Neutron Position Scenario]
  \label{def:physics:phyx-mini-0503:phyxminiproblems-problemphyxmini0503-matchesneutronpositionscenario}
  \lean{PhyXMiniProblems.ProblemPhyXMini0503.MatchesNeutronPositionScenario}
  \uses{def:physics:phyx-mini-0503:phyxminiproblems-problemphyxmini0503-onedimensionalquantumpositionsetup}
  The prose specifies a neutron and the graph specifies femtometer coordinates
\end{definition}
\begin{proof}
  This is the declared model definition of Matches Neutron Position Scenario.
\end{proof}
\begin{definition}[Matches Problem And Figure Readouts]
  \label{def:physics:phyx-mini-0503:phyxminiproblems-problemphyxmini0503-matchesproblemandfigurereadouts}
  \lean{PhyXMiniProblems.ProblemPhyXMini0503.MatchesProblemAndFigureReadouts}
  \uses{def:physics:phyx-mini-0503:phyxminiproblems-problemphyxmini0503-onedimensionalquantumpositionsetup}
  Numerical and textual data read directly from the question and raster. The endpoint 4, query cutoff 2, axis ticks, and printed labels occur here, but the requested event probability does not
\end{definition}
\begin{proof}
  This is the declared model definition of Matches Problem And Figure Readouts.
\end{proof}
\begin{definition}[Matches Supplied Probability Density Figure]
  \label{def:physics:phyx-mini-0503:phyxminiproblems-problemphyxmini0503-matchessuppliedprobabilitydensityfigure}
  \lean{PhyXMiniProblems.ProblemPhyXMini0503.MatchesSuppliedProbabilityDensityFigure}
  \uses{def:physics:phyx-mini-0503:phyxminiproblems-problemphyxmini0503-onedimensionalquantumpositionsetup, def:physics:phyx-mini-0503:phyxminiproblems-problemphyxmini0503-positiondensityreadout}
  The supplied V-shaped graph, transcribed as a piecewise-linear density. Inside the endpoints the height is a |x| / supportHalfWidth; outside it is zero. The first field identifies the plotted curve with this setup's wave-function density rather than with an unrelated scalar function
\end{definition}
\begin{proof}
  This is the declared model definition of Matches Supplied Probability Density Figure.
\end{proof}
\begin{definition}[Has Physical Probability Density Parameters]
  \label{def:physics:phyx-mini-0503:phyxminiproblems-problemphyxmini0503-hasphysicalprobabilitydensityparameters}
  \lean{PhyXMiniProblems.ProblemPhyXMini0503.HasPhysicalProbabilityDensityParameters}
  \uses{def:physics:phyx-mini-0503:phyxminiproblems-problemphyxmini0503-onedimensionalquantumpositionsetup}
  Positivity and ordering conditions implicit in a physical density graph
\end{definition}
\begin{proof}
  This is the declared model definition of Has Physical Probability Density Parameters.
\end{proof}
\begin{definition}[Satisfies One Dimensional Born Rule]
  \label{def:physics:phyx-mini-0503:phyxminiproblems-problemphyxmini0503-satisfiesonedimensionalbornrule}
  \lean{PhyXMiniProblems.ProblemPhyXMini0503.SatisfiesOneDimensionalBornRule}
  \uses{def:physics:phyx-mini-0503:phyxminiproblems-problemphyxmini0503-onedimensionalquantumpositionsetup, def:physics:phyx-mini-0503:phyxminiproblems-problemphyxmini0503-positiondensityreadout}
  The one-dimensional Born rule for arbitrary measurable position events. This is a governing law, not a statement of the probability of the particular event |x| ≥ 2 fm
\end{definition}
\begin{proof}
  This is the declared model definition of Satisfies One Dimensional Born Rule.
\end{proof}
\begin{definition}[Answer Choice]
  \label{def:physics:phyx-mini-0503:phyxminiproblems-problemphyxmini0503-answerchoice}
  \lean{PhyXMiniProblems.ProblemPhyXMini0503.AnswerChoice}
  Labels of the four probabilities displayed with the problem
\end{definition}
\begin{proof}
  This is the declared model definition of Answer Choice.
\end{proof}
\begin{definition}[displayed Probability]
  \label{def:physics:phyx-mini-0503:phyxminiproblems-problemphyxmini0503-displayedprobability}
  \lean{PhyXMiniProblems.ProblemPhyXMini0503.displayedProbability}
  \uses{def:physics:phyx-mini-0503:phyxminiproblems-problemphyxmini0503-answerchoice}
  Dimensionless probability printed beside each answer label
\end{definition}
\begin{proof}
  This is the declared model definition of displayed Probability.
\end{proof}
\begin{definition}[Matches Displayed Probability]
  \label{def:physics:phyx-mini-0503:phyxminiproblems-problemphyxmini0503-matchesdisplayedprobability}
  \lean{PhyXMiniProblems.ProblemPhyXMini0503.MatchesDisplayedProbability}
  \uses{def:physics:phyx-mini-0503:phyxminiproblems-problemphyxmini0503-onedimensionalquantumpositionsetup, def:physics:phyx-mini-0503:phyxminiproblems-problemphyxmini0503-requestedpositionprobability, def:physics:phyx-mini-0503:phyxminiproblems-problemphyxmini0503-answerchoice, def:physics:phyx-mini-0503:phyxminiproblems-problemphyxmini0503-displayedprobability}
  A displayed choice agrees exactly with the modeled event probability
\end{definition}
\begin{proof}
  This is the declared model definition of Matches Displayed Probability.
\end{proof}
\begin{definition}[Is Unique Matching Answer Choice]
  \label{def:physics:phyx-mini-0503:phyxminiproblems-problemphyxmini0503-isuniquematchinganswerchoice}
  \lean{PhyXMiniProblems.ProblemPhyXMini0503.IsUniqueMatchingAnswerChoice}
  \uses{def:physics:phyx-mini-0503:phyxminiproblems-problemphyxmini0503-onedimensionalquantumpositionsetup, def:physics:phyx-mini-0503:phyxminiproblems-problemphyxmini0503-answerchoice, def:physics:phyx-mini-0503:phyxminiproblems-problemphyxmini0503-matchesdisplayedprobability}
  A choice is the unique displayed probability agreeing with the model
\end{definition}
\begin{proof}
  This is the declared model definition of Is Unique Matching Answer Choice.
\end{proof}
\begin{lemma}[peak Density Readout eq one fourth]
  \label{lem:physics:phyx-mini-0503:phyxminiproblems-problemphyxmini0503-peakdensityreadout-eq-one-fourth}
  \lean{PhyXMiniProblems.ProblemPhyXMini0503.peakDensityReadout_eq_one_fourth}
  \uses{def:physics:phyx-mini-0503:phyxminiproblems-problemphyxmini0503-onedimensionalquantumpositionsetup, def:physics:phyx-mini-0503:phyxminiproblems-problemphyxmini0503-matchesproblemandfigurereadouts, def:physics:phyx-mini-0503:phyxminiproblems-problemphyxmini0503-matchessuppliedprobabilitydensityfigure, def:physics:phyx-mini-0503:phyxminiproblems-problemphyxmini0503-hasphysicalprobabilitydensityparameters, def:physics:phyx-mini-0503:phyxminiproblems-problemphyxmini0503-satisfiesonedimensionalbornrule}
  Normalization, the Born rule, and the two congruent triangular halves of the graph force the endpoint label a to read 1/4 fm⁻¹. This is a derived intermediate result, not a figure premise
\end{lemma}
\begin{proof}
  The conclusion follows from the governing relations and figure readouts named in the statement of peak Density Readout eq one fourth.
\end{proof}
% --- Archon declaration topology end ---
% --- Archon physics formalization source end ---
